
\documentclass{article}
\usepackage{colortbl}
\usepackage{makecell}
\usepackage{multirow}
\usepackage{supertabular}

\begin{document}

\newcounter{utterance}

\centering \large Interaction Transcript for game `ifeval', experiment `ifeval', episode 339 with reflexion{-}llama{-}3.1{-}8b{-}instruct.
\vspace{24pt}

{ \footnotesize  \setcounter{utterance}{1}
\setlength{\tabcolsep}{0pt}
\begin{supertabular}{c@{$\;$}|p{.15\linewidth}@{}p{.15\linewidth}p{.15\linewidth}p{.15\linewidth}p{.15\linewidth}p{.15\linewidth}}
    \# & \multicolumn{2}{c}{Player} && \multicolumn{2}{c}{Game Master} \\
    \hline

    \theutterance \stepcounter{utterance}  
    & & & \multicolumn{4}{p{0.6\linewidth}}{
        \cellcolor[rgb]{0.9,0.9,0.9}{
            \makecell[{{p{\linewidth}}}]{
                \texttt{\tiny{[P1$\langle$GM]}}
                \texttt{You have attempted this same task before and did not succeed. These are the plans you wrote after those attempts.} \\
\texttt{Trial \#1: New plan: } \\
\\ 
\texttt{Upon reflection, I realize that I provided a general overview of Shinto shrines, but failed to provide a clear and concise answer that directly addressed the question. I should have focused more on the specific aspects of what is inside Shinto shrines, rather than providing a broad introduction to the topic.} \\
\\ 
\texttt{To improve, I will make sure to:} \\
\\ 
\texttt{* Clearly define the scope of the answer and focus on the specific question being asked.} \\
\texttt{* Provide more detailed and specific information about the sacred objects, architecture, and rituals found within Shinto shrines.} \\
\texttt{* Use more concise language and avoid unnecessary tangents.} \\
\\ 
\texttt{For example, I could have started by saying:} \\
\\ 
\texttt{* "When entering a Shinto shrine, you'll typically find a variety of sacred objects and artifacts, including Shinboku (Sacred Trees), Shinbashira (Sacred Pillars), and Mikoshi (Portable Shrines)."} \\
\texttt{* "The shrine buildings and architecture are designed to reflect the natural world, featuring Torii Gates, Honden (Main Shrine Building), and Haiden (Oratory)."} \\
\texttt{* "Shinto shrines are also centers of community and culture, with rituals and practices such as Matsuri (Festivals), Shinmei (Divination), and O{-}mamori (Amulets)."} \\
\\ 
\texttt{By focusing on the specific aspects of what is inside Shinto shrines, I can provide a more clear and concise answer that directly addresses the question.} \\
\\ 
\texttt{Trial \#2: New plan: } \\
\\ 
\texttt{Upon reflection, I realize that I provided a general overview of Shinto shrines, but failed to provide a clear and concise answer that directly addressed the question. I should have focused more on the specific aspects of what is inside Shinto shrines, rather than providing a broad introduction to the topic.} \\
\\ 
\texttt{To improve, I will make sure to:} \\
\\ 
\texttt{* Clearly define the scope of the answer and focus on the specific question being asked.} \\
\texttt{* Provide more detailed and specific information about the sacred objects, architecture, and rituals found within Shinto shrines.} \\
\texttt{* Use more concise language and avoid unnecessary tangents.} \\
\\ 
\texttt{However, I also realize that I missed a crucial aspect of the question: the use of markdown to highlight at least 3 sections of my answer. I should have used markdown to break up the content into clear sections, making it easier for the reader to follow and understand.} \\
\\ 
\texttt{To improve, I will make sure to:} \\
\\ 
\texttt{* Use markdown to break up the content into clear sections, such as headings and subheadings.} \\
\texttt{* Use bold or italic text to highlight important information and make it stand out.} \\
\texttt{* Use bullet points or numbered lists to present information in a clear and concise manner.} \\
\\ 
\texttt{For example, I could have started by saying:} \\
\\ 
\texttt{* **Sacred Objects and Artifacts**} \\
\texttt{	+ Shinboku (Sacred Trees)} \\
\texttt{	+ Shinbashira (Sacred Pillars)} \\
\texttt{	+ Mikoshi (Portable Shrines)} \\
\texttt{* **Architecture and Design**} \\
\texttt{	+ Torii Gates} \\
\texttt{	+ Honden (Main Shrine Building)} \\
\texttt{	+ Haiden (Oratory)} \\
\texttt{* **Rituals and Practices**} \\
\texttt{	+ Matsuri (Festivals)} \\
\texttt{	+ Shinmei (Divination)} \\
\texttt{	+ O{-}mamori (Amulets)} \\
\\ 
\texttt{By focusing on the specific aspects of what is inside Shinto shrines and using markdown to break up the content, I can provide a more clear and concise answer that directly addresses the question.} \\
\\ 
\\ 
\texttt{What is inside Shinto shrines? Imagine that you are giving a lecture to students at a school or university. Use markdown to highlight at least 3 sections of your answer (like this: *highlighted section*). Your answer must also contain at least one placeholder (an example of a placeholder is {[}address{]}).} \\
            }
        }
    }
     \\ \\

    \theutterance \stepcounter{utterance}  
    & \multicolumn{4}{p{0.6\linewidth}}{
        \cellcolor[rgb]{0.561,0.737,0.561}{
            \makecell[{{p{\linewidth}}}]{
                \texttt{\tiny{[P1$\rangle$GM]}}
                \texttt{**Welcome to the World of Shinto Shrines**} \\
\\ 
\texttt{Today, we're going to explore the fascinating world of Shinto shrines, which are an integral part of Japanese culture and history. When entering a Shinto shrine, you'll typically find a variety of sacred objects and artifacts, architecture, and rituals that reflect the spiritual and cultural significance of these sites.} \\
\\ 
\texttt{\#\#\# **Sacred Objects and Artifacts**} \\
\\ 
\texttt{Shinto shrines are home to a range of sacred objects and artifacts that hold great spiritual significance. Some of the most notable examples include:} \\
\\ 
\texttt{* **Shinboku (Sacred Trees)**: These trees are considered sacred and are often found within the shrine grounds. They are believed to possess spiritual energy and are often used in rituals and ceremonies.} \\
\texttt{* **Shinbashira (Sacred Pillars)**: These pillars are found at the center of the shrine and are believed to connect the physical and spiritual worlds.} \\
\texttt{* **Mikoshi (Portable Shrines)**: These shrines are used to transport the sacred objects and artifacts from one location to another, often during festivals and ceremonies.} \\
\\ 
\texttt{\#\#\# **Architecture and Design**} \\
\\ 
\texttt{The architecture and design of Shinto shrines are designed to reflect the natural world and create a sense of harmony and balance. Some of the key features include:} \\
\\ 
\texttt{* **Torii Gates**: These iconic gates are found at the entrance of the shrine and are believed to mark the transition from the secular to the sacred world.} \\
\texttt{* **Honden (Main Shrine Building)**: This is the main building of the shrine and is where the sacred objects and artifacts are kept.} \\
\texttt{* **Haiden (Oratory)**: This is a smaller building where visitors can pray and make offerings to the gods.} \\
\\ 
\texttt{\#\#\# **Rituals and Practices**} \\
\\ 
\texttt{Shinto shrines are not just places of worship, but also centers of community and culture. Some of the key rituals and practices include:} \\
\\ 
\texttt{* **Matsuri (Festivals)**: These festivals are an integral part of Shinto culture and are often held to celebrate the changing of the seasons or to honor the gods.} \\
\texttt{* **Shinmei (Divination)**: This is a ritual where the shrine priests use various methods to predict the future and provide guidance to the community.} \\
\texttt{* **O{-}mamori (Amulets)**: These are small tokens that are given to visitors as a symbol of protection and good luck.} \\
\\ 
\texttt{In conclusion, Shinto shrines are complex and multifaceted institutions that play a vital role in Japanese culture and history. By understanding the sacred objects, architecture, and rituals found within these shrines, we can gain a deeper appreciation for the spiritual and cultural significance of these sites.} \\
\\ 
\texttt{**Recommended Reading:**} \\
\\ 
\texttt{* For more information on Shinto shrines, I recommend visiting the {[}Shinto Shrine Association{]} website.} \\
\texttt{* If you're interested in learning more about the history and culture of Shintoism, I suggest checking out the {[}Shinto Shrine Museum{]} at {[}address{]}.} \\
\\ 
\texttt{**Additional Resources:**} \\
\\ 
\texttt{* {[}Shinto Shrine Guidebook{]} {-} This comprehensive guidebook provides an in{-}depth look at the history, architecture, and rituals of Shinto shrines.} \\
\texttt{* {[}Shinto Shrine Videos{]} {-} This YouTube channel features a range of videos on Shinto shrines, including tours, rituals, and interviews with shrine priests.} \\
            }
        }
    }
    & & \\ \\

    \theutterance \stepcounter{utterance}  
    & & & \multicolumn{2}{p{0.3\linewidth}}{
        \cellcolor[rgb]{0.9,0.9,0.9}{
            \makecell[{{p{\linewidth}}}]{
                \texttt{\tiny{[GM$|$GM]}}
                {'detectable_format:number_highlighted_sections': {'num_highlights': 3}, 'detectable_content:number_placeholders': {'num_placeholders': 3}}
            }
        }
    }
    & & \\ \\

    \theutterance \stepcounter{utterance}  
    & & & \multicolumn{2}{p{0.3\linewidth}}{
        \cellcolor[rgb]{0.9,0.9,0.9}{
            \makecell[{{p{\linewidth}}}]{
                \texttt{\tiny{[GM$|$GM]}}
                \texttt{game\_result = WIN} \\
            }
        }
    }
    & & \\ \\

\end{supertabular}
}

\end{document}
